% This part was made obsolete on 2024-10-30.
% I guess that it is not necessary to summarise my book here.

\documentclass[../libro]{subfiles}

\begin{document}

\ifSubfilesClassLoaded{\appendix\chapter{后日谈}\clearpage}{}

\section{总结}

用归纳定义定义行列式的线性代数教材%
似乎比用组合定义定义行列式的线性代数教材少.
具体地, 写给数学学生的教材少地用归纳定义;
一些写给非数学学生的教材用归纳定义,
但不多.
我在本书用归纳定义,
定义了行列式,
并用 \(\rho\)-记号简单地证明了行列式的公式
(如按一列展开行列式,
按多列展开行列式,
完全展开行列式)
与性质,
从而说明了,
归纳定义也是一个好的定义.

我用行列式讲了线性方程组的理论.
我想, 这或许是行列式的一个好的应用.

我用行列式讲了 \(2\)~元 \(2\)~次式的理论,
与如何用它解 \(2\)~元 \(2\)~次方程组%
与 \(1\)~元 \(4\)~次方程.
我想, 这或许也是行列式的一个好的应用.

在本章 (本附录), 我讲了一些关于行列式的结论,
包括但不限于%
其他的用行列式的性质确定行列式的方式,
判断 (复) 阵的行列式不是 \(0\) 的方法,
判断实阵的行列式大于 \(0\) 的方法,
古伴的性质,
反称阵的行列式与 pfaffian 的关系,
辛阵的行列式为 \(1\).

我似乎讲完了我想讲的.
再见.

% 按多列 (行) 展开行列式是教学的一个难点:
% 我们少地用它算行列式.
% 我认为,
% 这是一个更理论的 (而不是应用的) 定理.
% 我想, 我要么不讲它, 要么通俗地讲它.
% 所以, 我加了一个具体的例,
% 并希望此例可助学生理解按多列展开行列式.
% 证明的方法还是数学归纳法.
% 我用到了按 (任何) 一列展开行列式的公式.
% 所以, 理论地,
% 这个论证可被任何一本教材用
% (一般地, 讲按多列展开行列式前,
% 都会讲按一列展开行列式).
% 不过, 我没看到这样的教材
% (包括一本用归纳定义定义行列式,
% 但用完全展开的公式证明此事%
% 的线性代数教材).

\end{document}

This is just a closing remark.
I summarise what topics that are probably interesting
are discussed in the book.
